\begin{abstract}
This study isolates threshold calibration scope as the sole experimental variable in a federated IoT anomaly-detection pipeline: a fixed FedAvg autoencoder (E=1, full participation) is trained once per seed, and the same encoder, seeds, and per-client score artifacts are reused across shared (B1), per-client (B2), family-mean (B3), and cluster-mean (B4) threshold policies. The primary comparison is B1 vs.\ B2 on N-BaIoT (K=9 physical devices, q=0.95, five seeds). Because B2 sets each client's threshold at its own benign 95th percentile, FPR equalization is expected by construction under stable calibration/test distributions; the empirical contribution is the held-out magnitude, seed consistency, detection-quality tradeoff, and B4's taxonomy-free approximation. Across the five seeds, B2 reduces CV(FPR) from 1.017 to 0.299 (delta 0.718, bootstrap CI [0.647, 0.769], all five seed deltas positive), and B4 recovers about 52\% without a taxonomy. The gain is not uniform: lower-tail Macro-F1 degradation concentrates in a single low-separability client, and under the randomized file-level CICIoT2023 partition (Regime~B; pseudo-clients, not physical devices), where benign distributions are near-homogeneous, no improvement is observed. These findings are bounded by nine devices, E=1 (which limits local drift), five seeds, and CICIoT2023 conclusions that hold only under the tested partition.
\end{abstract}

\begin{IEEEkeywords}
federated learning, IoT anomaly detection, threshold calibration, false-positive-rate dispersion, autoencoder, non-IID data, personalization.
\end{IEEEkeywords}
